% !TEX program = xelatex
% 编译方式:在 Overleaf 中将编译器设置为 XeLaTeX
\documentclass[11pt,a4paper]{article}

% ============ 宏包 ============
\usepackage[UTF8]{ctex}               % 中文支持
\usepackage[left=1.8cm,right=1.8cm,top=1.8cm,bottom=1.8cm]{geometry}
\usepackage{xcolor}
\usepackage{graphicx}
\usepackage{enumitem}
\usepackage{titlesec}
\usepackage{tikz}
\usepackage{fontawesome5}             % 图标(可选)
\usepackage{hyperref}
\usepackage{array}
\usepackage{tabularx}

% ============ 颜色定义 ============
\definecolor{themecolor}{RGB}{196, 30, 58}   % 红色主题
\definecolor{tagcolor}{RGB}{196, 30, 58}
\definecolor{textgray}{RGB}{80, 80, 80}

% ============ 超链接样式 ============
\hypersetup{
    colorlinks=true,
    urlcolor=themecolor,
    linkcolor=themecolor
}

% ============ 段落设置 ============
\setlength{\parindent}{0pt}
\pagestyle{empty}

% ============ 章节标题样式 ============
\titleformat{\section}
    {\color{white}\bfseries\large}
    {}
    {0pt}
    {\colorbox{themecolor}{\parbox{\dimexpr\textwidth-2\fboxsep}{\strut #1}}}
\titlespacing*{\section}{0pt}{12pt}{6pt}

% ============ 小标签(C9、985) ============
\newcommand{\tag}[1]{%
    \tikz[baseline=(X.base)]\node[rectangle, rounded corners=2pt,
    fill=tagcolor, text=white, inner sep=2pt, font=\footnotesize\bfseries] (X) {#1};%
}

% ============ 条目标题(左标题 + 右日期) ============
\newcommand{\entry}[2]{%
    \noindent\textbf{#1}\hfill{\color{textgray}#2}\\
}

\newcommand{\subentry}[2]{%
    \noindent#1\hfill{\color{textgray}#2}\\
}

% ============ 列表样式 ============
\setlist[itemize]{leftmargin=1.2em, itemsep=2pt, topsep=2pt, parsep=0pt}

% ============ 开始文档 ============
\begin{document}

% ====================== 个人信息 ======================
\begin{minipage}[t]{0.75\textwidth}
    \vspace{0pt}
    {\Huge\bfseries\color{themecolor} 李金美}\\[8pt]
    \faPhone\ 18790566306 \quad \faEnvelope\ 18790566306@163.com\\[4pt]
    23岁 \ | \ 女 \ | \ 籍贯:河南新乡 \ | \ 汉族 \ | \ 中共党员
\end{minipage}%
\hfill
\begin{minipage}[t]{0.22\textwidth}
    \vspace{0pt}
    \begin{flushright}
    \framebox[3cm]{\rule{0pt}{3.6cm}\hspace{-3cm}\parbox{3cm}{\centering \small 照片\\(可替换)}}
    \end{flushright}
\end{minipage}

\vspace{6pt}

% ====================== 教育背景 ======================
\section{教育背景}

\entry{中国科学技术大学 \tag{C9} \tag{985}}{2025年09月 - 至今}
\subentry{网络空间安全专业 \ 硕士（研一） \ 网络空间安全学院}{合肥}
\begin{itemize}
    \item \textbf{保研学硕,中国科学技术大学硕士一等学业奖学金}
\end{itemize}

\vspace{4pt}

\entry{武汉大学 \tag{985}}{2021年09月 - 2025年06月}
\subentry{网络空间安全专业 \ 本科 \ 国家网络安全学院}{武汉}
\begin{itemize}
    \item \textbf{GPA:3.92/4.0(前3\%)}
    \item \textbf{主修课程:}编译原理(93)、计算机网络(97)、计算机组成原理(95)、数据结构(91)、操作系统原理及安全(93)、密码学(95)、数据库系统及安全(95)、汇编语言及逆向分析(91)等
    \item \textbf{奖项荣誉:}国家奖学金、武汉大学优秀本科毕业生、雷军计算机创新与发展基金、武汉大学三好学生、武汉大学优秀学生奖学金
    \item \textbf{科研成果:}发明专利《基于声学传感的智能用户认证方法、设备及软件产品》第一作者,授权专利号:ZL 202410993112.5;2024年第十七届全国大学生信息安全竞赛(A类学科竞赛)国家级一等奖;2024年第十七届中国大学生计算机设计大赛(A类学科竞赛)中南地区赛一等奖
\end{itemize}

% ====================== 实习经历 ======================
\section{实习经历}

\entry{科大讯飞股份有限公司 ---- 讯飞医疗大模型实习}{2025年11月 - 2026年05月}

% -------- 项目一：病历质控 GRPO --------
\textbf{病历质控规则奖励 GRPO 训练（Qwen3.5-27B / Megatron / vLLM）}
\begin{itemize}
    \item \textbf{数据构造与难度配比:}基于线上 bad case 对模型思维链进行归因分析，抽样 500 条 CoT、15{,}249 个推理步骤，抽象出涵盖推理正确性、规则逻辑、知识准确性等 10 维 step-level 错误标签，定位错误主要集中在前 3 步；筛选难度适中（模型单步成功率 30\%--70\%）样本构建训练集，确保 GRPO 组内 reward 方差不为零，避免全 0/全 1 group 产生无效梯度。
    \item \textbf{规则奖励函数设计:}实现 \texttt{compute\_score\_blzk\_rule}，核心设计包括：(1) 自动剥离 Qwen3.5 native thinking 模式的 \texttt{<think>} 块（覆盖 chat template 追加开始标签缺失的边界情况）；(2) 支持 5 种合法 JSON key 组合的精确模式匹配，对格式变化鲁棒；(3) 全角引号等中文噪声规范化；产出 0/1 二元奖励并附加 \texttt{judge\_reason} 分类（ok / mismatch / json\_not\_found / json\_parse\_failed / key\_mode\_not\_matched），用于训练全程实时监控。
    \item \textbf{GRPO 训练配置:}基于 verl 框架，Megatron 后端（TP=4、PP=1）+ vLLM 异步 rollout（TP=8），采用 dapo reward manager 配合 overlong reward shaping 抑制长 response 截断引入的奖励噪声。
    \item \textbf{训练效果:}在 2{,}000 条内部病历质控测试集上，F1 值从基线 \textbf{87.4\%} 提升至 \textbf{89.27\%}；观察到精确率上升而召回率下降，分析为 RL 使模型在不确定时趋向保守（倾向保留"合格"判断规避错误），计划引入更多难负样本调整难度分布以修正。
    \item \textbf{熵坍塌防御实验（设计中）:}观察到训练中 entropy 呈下降趋势；设计两组对照实验：(A) 开启 KL loss（$\lambda_\text{KL}=0.01$）vs 关闭；(B) 对称 clip（$\varepsilon=0.2$）vs 非对称 clip-higher（$\varepsilon_\text{high}=0.28$），分析各配置对 entropy 维持与最终 F1 的影响。
\end{itemize}

% -------- 项目二：病历生成 GenRM（进行中）--------
\textbf{病历生成 GenRM GRPO 训练（进行中）}
\begin{itemize}
    \item \textbf{任务背景与难点:}从医患对话自动生成结构化病历，属于开放式生成任务，无标准参考答案；Qwen3.5-27B zero-shot 与 SFT 版本在 LLM-as-Judge 评测下得分均不理想；规则奖励无法处理生成质量的语义维度，选用生成式奖励模型（GenRM）驱动 GRPO。
    \item \textbf{GenRM 方案设计:}以大模型作为奖励模型，对生成病历的完整性、结构规范性、信息忠实度等维度进行细粒度评分；设计对应评测 prompt 与 rubric，确保奖励信号与人工评估高度对齐。
    \item \textbf{数据与评测:}整理医患对话--病历配对数据，通过 LLM 打分建立定量基线；\textbf{[待补充:数据量、评测基线分数]}。
\end{itemize}

% -------- 项目三（注释掉）：医考选择题 GenRM --------
% \textbf{医考选择题生成式奖励模型 GRPO 训练（进行中）}
% \begin{itemize}
%     \item \textbf{数据构造:}整理构造 12{,}950 条中文多选题与 65{,}293 条中英文单选题,将单选题筛选并改写为开放式问答,
%           最终保留 53{,}668 条高质量医考开放式 QA;组织多模型采样回答并用 LLM 自动打分,
%           构建训练集 20{,}000 query × 8 answers = 160{,}000 条偏好数据候选集与 4{,}000 条验证集。
%     \item \textbf{生成式奖励模型(GenRM)设计:}采用 Qwen2.5-7B-Instruct 作为 GenRM 后端。
%     \item \textbf{对比实验设计(计划):}规则匹配 baseline / 判别式 RM / 生成式 RM 三组对比。
% \end{itemize}

% -------- 项目四：环境搭建 --------
\textbf{Qwen3.5-27B 多机 GRPO + SFT 训练环境从零搭建}
\begin{itemize}
    \item \textbf{背景:}原业务依托华为昇腾集群 + 公司自研训练框架 + 自研基模；因自研基模效果不及预期，决策切换至 Qwen3.5-27B，但无现成环境，由本人从零完成 H200 多机集群 post-training 全套工具链部署（verl / Megatron-LM / vLLM / Ray）。
    \item \textbf{关键工程问题:}解决多机部署中\textbf{多网卡 IP 解析错误、Ray 节点资源 race condition、shell 解析 Ray INFO 日志污染、Hydra struct mode 兼容、FlashInfer JIT 在 NFS 上并发竞争}等问题；针对 Qwen3.5 GDN 线性注意力不支持 THD packed sequence 的架构限制，确定 \texttt{use\_remove\_padding=False} + \texttt{use\_dynamic\_bsz=False} 的 BSHD 路径方案。
    \item \textbf{SFT 环境与分析:}配置 Swift + H200 + Qwen3.5-27B 全量 SFT 环境；对 230 万条医学数据（48卡、约3天）SFT 后指令遵循能力（IFEval）下降问题进行分析，定位为思维链显著缩短、训练数据原为自研基模定制导致分布不对齐，提供数据质量改进建议。
    \item \textbf{工程方法论:}梳理业务侧快速迭代模式：先以单业务小批量数据验证数据有效性，再合入全业务大版本（230 万条）进行正式训练；掌握三层 batch 设计、KL 加 loss vs 加 reward、低方差 KL 估计器（k3）、PPO clip-higher、reward manager 选型等核心配置与权衡。
\end{itemize}

% -------- 项目五：多模态评测 --------
\textbf{医疗多模态大语言模型自动评测与错误分析}
\begin{itemize}
    \item 设计覆盖多维度理解、逻辑推理、医学正确性、安全性、回答有用性、格式规范等维度的自动评测框架，辅助大模型迭代中的回归与对比分析。
    \item 基于 LLM-as-a-Judge 设计并优化评测 prompt、rubric 细化与子分标签；调用 Claude 等模型进行结构化打分，完成维度分析、bad case 回溯与评测报告可视化。
    \item 针对医学场景常见错误类型（笼统回答、幻觉、证据缺失、图像利用不足、结构输出不规范等）构建错误分类视角，为模型迭代与 benchmark 设计提供依据。
    \item \textbf{多模态 SFT 跟进:}参与多模态 SFT 数据梳理与分析（Swift + H200 + Qwen3.5-27B，冻结视觉编码器，约 2 万图文 QA 数据，16 卡 batchsize=16，约 10 小时），了解各业务数据分布（医考 / 诊疗助理 / 多轮对话等）及数据有效性验证流程。
\end{itemize}

% -------- 项目六：KIE Benchmark --------
\textbf{医疗 KIE Benchmark 构建}
\begin{itemize}
    \item 参与医疗信息抽取场景下 benchmark 任务设计，覆盖分类、字段值抽取、关系判断与复杂结构输出等多种任务类型。
\end{itemize}

% ====================== 个人总结 ======================
\section{个人总结}

\textbf{充足的实习时间:}可实习6个月以上\\[4pt]
\textbf{技术技能}
\begin{itemize}
    \item \textbf{大模型训练框架:}verl（GRPO/PPO 全流程）、Swift（SFT）、Megatron-LM、vLLM、Transformers
    \item \textbf{大模型评测与分析:}LLM-as-a-Judge、奖励函数设计、Rubric 设计、Bad Case 分析、Post-training Evaluation
    \item \textbf{医疗 AI:}Medical QA、病历质控、病历生成、KIE、多模态评测、结构化信息抽取
    \item \textbf{工程与实验:}Python、PyTorch、Shell、Linux、多机 GPU 集群调试、Ray 集群编排
\end{itemize}

\end{document}
